\section{Norm uniqueness of sequence models}\label{sec:uniqueness}

We prove uniqueness for the models of \Cref{seq:model}, allowing any of
the spaces $Y$ and measures $m$ specified there.  The uniqueness
argument uses the coarse labels
\[
 S=S_Y=C_0(Y)\otimes P,\qquad S(U)=C_0(Y)\otimes P(U),
 \qquad U\subseteq X\text{ open}.
\]
The same statements apply if the prescribed measure is multiplied by a
strictly positive constant.  A coordinate unitary will always mean a
unitary in $\mathcal M(Q)$; convergence of coordinate conjugacies is in
norm on each fixed input.
\begin{theorem}[Norm uniqueness]\label{un:uniqueness}
Let $p,q:S_Y\to Q_\infty$ be models for the same data.  For any
c.p.c. lifts $p_n,q_n:S_Y\to Q$, there are
$u_n\in\mathcal U(\mathcal M(Q))$ such that
\[
 \lim_{n\to\infty}\|u_np_n(a)u_n^*-q_n(a)\|=0,
 \qquad a\in S_Y.
\]
\end{theorem}

There are two steps.  Equal trace data first give, after conjugacy and
rational averaging, a difference in every local error ideal $J(U)$.
A small fraction of a common spare model then absorbs that difference
in norm.  Splitting each original model into two UHF halves supplies
the spare: change the first half while keeping the second, and then
change the second while keeping the first.  The intermediate absorption
statements allow a general labeling space, so that the cone correction
can use them with its finer ideal labels.

\subsection{Compact sets of weights and tracial alignment}

\begin{lemma}[Controlled trace evaluations]\label{un:trace-cap}
Let $C$ be a $C^*$-algebra and $b\in(C\otimes\K)_+$.  The set
\[
 K_b=\{\gamma\in T(C):d_\gamma(b)\leq1\}
\]
is compact and convex.  If $x\in(C\otimes\K)_+$ satisfies
$\sup_{\gamma\in K_b}d_\gamma(x)\leq L<\infty$, then
$\gamma\mapsto\gamma(x)$ is finite, continuous, and affine on $K_b$.
The constant $L$ may depend on $x$.

Consequently this applies separately to all the first and square
evaluations of a fixed completely positive candidate whose tested
values have bounded rank on $K_b$.  It also applies to
\[
 \gamma\longmapsto
 \gamma\bigl(\lvert u x_1u^*-x_2\rvert^2\bigr),
 \qquad u\in\mathcal U(C^{\sim}),
\]
whenever the positive elements $x_1,x_2$ have uniformly bounded ranks
there, with a rank bound independent of $u$.
\end{lemma}

\begin{proof}
Rank evaluation is lower semicontinuous on the compact trace cone, so
$K_b$ is closed.  Its convexity follows from additivity and positive
homogeneity of rank.  For $\eta>0$ the assumed rank bound gives
\begin{equation}\label{un:cut-loss}
 0\leq\gamma(x)-\gamma((x-\eta)_+)\leq\eta L,
 \qquad \gamma\in K_b,
\end{equation}
and $\gamma(x)\leq\|x\|L$.  If $\gamma_i\to\gamma$ in $K_b$, the
defining cut inequalities for the trace topology imply
\[
 \gamma(x)\leq\liminf_i\gamma_i(x)
 \leq\limsup_i\gamma_i(x)
 \leq\gamma(x)+\eta L.
\]
Letting $\eta\downarrow0$ proves continuity.  Affineness follows from
the cone operations; it does not require cancellation in that cone.
Squaring a positive element does not change its support rank.  Finally,
support subadditivity and unitary invariance give
\[
 d_\gamma\bigl(\lvert u x_1u^*-x_2\rvert^2\bigr)
 \leq d_\gamma(x_1)+d_\gamma(x_2),
\]
which proves the last assertion.
\end{proof}

\begin{lemma}[Barycenters on a trace cap]\label{un:barycenter}
Let $K_b$ be as in \Cref{un:trace-cap}.  For every probability measure
$\mu$ on $K_b$ there is $\bar\gamma\in K_b$ such that
\begin{equation}\label{un:barycenter-formula}
 v(\bar\gamma)=\int_{K_b}v\,d\mu
\end{equation}
for every finite continuous affine real function $v$ on $K_b$.
For a positive measure of finite mass $M>0$, its integral is instead
$M v(\bar\gamma)$, with $\bar\gamma$ chosen for $\mu/M$; zero measures
contribute zero.  Signed measures may be treated by their two positive
parts.
\end{lemma}

\begin{proof}
Approximate $\mu$ weakly by finitely supported probability measures
$\mu_j=\sum_k t_{jk}\delta_{\gamma_{jk}}$, omitting terms of weight
zero.  Their finite barycenters
$\beta_j=\sum_k t_{jk}\gamma_{jk}$ belong to $K_b$.  Take a subnet
converging in the compact set $K_b$ to $\bar\gamma$.  For every $v$ in
the statement,
\[
 v(\bar\gamma)=\lim_jv(\beta_j)
 =\lim_j\int v\,d\mu_j=\int v\,d\mu.
\]
The same subnet works for all these functions.  Neither uniqueness of
$\bar\gamma$ on discontinuous evaluations nor subtraction of extended
weights is asserted or used.
\end{proof}

\Needspace{3\baselineskip}
\begin{lemma}[Alignment for finitely many weights]\label{un:finite-traces}
Let $C$ be separable, let $S$ be separable and nuclear, and let
$p,q:S\to C$ be homomorphisms with equal pullbacks on $T(C)$.
For a finite list of pairs $(s_i,\gamma_i)$, where $s_i\in S_+$ and
\[
 d_{\gamma_i}(p(s_i))+d_{\gamma_i}(q(s_i))<\infty,
\]
and any $\delta>0$, there is $u\in\mathcal U(M_2(C)^{\sim})$ such that
\begin{equation}\label{un:finite-L2}
 (\gamma_i\otimes\Tr_2)
 \bigl(\lvert u(p(s_i)\oplus0)u^*-(q(s_i)\oplus0)\rvert^2\bigr)<\delta
 \quad\text{for every }i.
\end{equation}
The added matrix direction is internal and has unnormalized trace.
\end{lemma}

\begin{proof}
The equality of pullback weights implies equality of all their ranks.
By \cite[Lemma~5.1]{CGNN}, the support projections of $p(s)$ and $q(s)$
are equivalent in $C^{**}$ for every $s\in S_+$.  Their images remain
equivalent under the normal extension of every representation of $C$.

The finite-rank assumption puts $p(s_i)$ and $q(s_i)$ in
$\Fin(\gamma_i)$. After zero padding, $M_2(\Fin(\gamma_i))$ is invariant under
$M_2(C)^{\sim}$-unitaries, so the conjugates and squared differences
below remain in that finite ideal.
For each selected weight, take its trace representation $\pi_i$ on
its densely finite ideal after quotienting by its zero ideal, and
extend it nondegenerately to $C$.  Let $\Tr_i$ be the resulting normal
semifinite trace on $M_i=\pi_i(C)''$.  Form
$\pi=\pi_0\oplus\pi_1\oplus\cdots\oplus\pi_r$, with $\pi_0$
a faithful representation on a separable Hilbert space, and put
$M=\pi(C)''$.  The trace representations can also be chosen separably
acting: their densely finite ideals are separable and have countable
approximate units of finite-trace cuts.  Thus $M$ has separable predual.
Restriction to the $i$th reducing Hilbert summand gives a normal
surjective homomorphism $q_i:M\to M_i$.  Its kernel is
$M(1-z_i)$ for a central projection $z_i$, and $q_i$ identifies
$Mz_i$ with $M_i$.  Define
\[
 \widehat\gamma_i=\Tr_i\circ q_i.
\]
This is the transported normal semifinite trace on $Mz_i$, extended
by zero on $M(1-z_i)$.  The projections $z_i$ may overlap; the
construction is made separately for each weight.  On the densely
finite ideal of $\gamma_i$ one has
$\widehat\gamma_i\pi=\gamma_i$, so it computes all designated
finite-rank tests and their squares.

We apply von Neumann uniqueness only after a common zero padding and
unitization.  Let $S^\dagger$ be the forced unitization, and define
unital homomorphisms into $M_2(M)$ by
\[
 \widehat p(a+\lambda1)=(\pi p(a)\oplus0)+\lambda1_{M_2(M)},\qquad
 \widehat q(a+\lambda1)=(\pi q(a)\oplus0)+\lambda1_{M_2(M)}.
\]
These maps have equal support ranks on every positive element of
$S^\dagger$.  If its scalar part is zero, this follows from the
support equivalences already proved.  If the scalar part is positive,
then on the properly infinite central part of $M$ both support
projections contain the second matrix corner, which is equivalent to
the unit of $M_2(M)$ there.  On the finite central part, every normal
finite trace has equal pullbacks along $\pi p$ and $\pi q$; this follows
by restricting it to $C$ and using the hypothesis.  The unital
extensions also have equal pullbacks, hence equal ranks by monotone
spectral calculus.  Normal traces determine equivalence of projections
in a finite von Neumann algebra.  This proves the claim on both parts.

The source $S^\dagger$ is nuclear.  We use the unital form of von Neumann
uniqueness in \cite[Proposition~2.1, Corollary~3.10, and Section~4]{CGNN}.
For completeness, the properly infinite factor argument needs only the
partial-isometry approximation of \cite[Lemma~3.6]{CGNN} and the following
unital completion.  Including $1$ among the inputs gives partial
isometries $v$ whose initial projections $e=v^*v$ tend strongly to $1$.
For finitely many normal states choose $e_0\leq e$ arbitrarily close to
$e$ on those states, with
$1-e_0\sim1-ve_0v^*\sim1$: take $e_0=e$ if $e$ is finite, and otherwise
remove an infinite tail decreasing strongly to zero.  Complete $ve_0$
to a unitary $u$.  Since $(u-v)e_0=0$, Cauchy--Schwarz gives
\[
 |\rho(u^*xu-v^*xv)|\leq4\|x\|\rho(1-e_0)^{1/2}.
\]
Thus the partial-isometry conjugacies can be replaced by unitary
conjugacies in the weak star topology; \cite[Lemma~3.9]{CGNN} makes
them strong star.  Section~4 then assembles these unital factor conclusions by
central decomposition; separability lets one impose the support
comparisons on one countable dense set of positive cuts before
selecting the factorwise unitaries.  We obtain $w_j\in M_2(M)$ with
\[
 w_j(\pi p(a)\oplus0)w_j^*\longrightarrow\pi q(a)\oplus0
 \quad\text{strong star},\qquad a\in S.
\]
The unital extensions are needed here: support ranks of nonunital maps
alone do not compare the complements of their support projections.

For clarity, the passage to $L^2$ uses more than strong convergence.
Put $x_j=w_j(\pi p(s_i)\oplus0)w_j^*$ and
$x=\pi q(s_i)\oplus0$.  Write
$T_i=\widehat\gamma_i\otimes\Tr_2$.  Their $T_i$-norms are finite and satisfy
\[
 \|x_j\|_{2,T_i}^2
 =\gamma_i(p(s_i^2))
 =\gamma_i(q(s_i^2))
 =\|x\|_{2,T_i}^2.
\]
For a bounded test $y\in L^1(T_i)\cap
L^2(T_i)$ supported on trace-finite projections, its
pairing with $x_j-x$ is the normal bounded functional
$t\mapsto T_i(y^*t)$ and tends to zero.  Such tests
are dense in $L^2(T_i)$; the common
$L^2$ bound therefore gives weak $L^2$ convergence.  Equality of norms
now gives $\|x_j-x\|_{2,T_i}\to0$.

Choose one $w_j$ meeting the finite list to sufficiently small
tolerance.  A bounded self-adjoint logarithm of this unitary can be
approximated strong star, within its norm bound, by self-adjoint
elements of $M_2(\pi(C))^{\sim}$, by Kaplansky density
\cite[Theorem~1]{KaplanskyDensity}.  Exponentiation
gives unitaries from $M_2(C)^{\sim}$ converging strong star to $w_j$.
Their conjugates again have the same designated $L^2$ norms, so the
preceding finite-trace-test argument gives $L^2$ convergence for this
approximation too.  On these conjugates the normal traces compute
the original weights, since their finite ideals are invariant under
$M_2(C)^{\sim}$-unitaries.  This proves \eqref{un:finite-L2}.
\end{proof}

\begin{lemma}[Tracial alignment]\label{un:alignment}
If $p,q:S\to D$ are models with the same data, then rational diagonal
averaging, coordinate multiplier conjugacy, and the foldings of
\Cref{st:folding} give models $\widetilde p,\widetilde q:S\to D$ with
\begin{equation}\label{un:local-difference}
 (\widetilde p-\widetilde q)(S(U))\subseteq J(U)
 \qquad(U\subseteq X\text{ open}).
\end{equation}
Each new model is coordinate approximately unitarily equivalent to
its respective original model.  The operations may be performed in
one changing block while a prescribed spare block is carried along
unchanged, apart from its common amplification and folding.
\end{lemma}

\begin{proof}
By \Cref{seq:all-traces}, choose a separable subalgebra $D_0$ containing
the two ranges, the required constant representatives, and the
comparison witnesses, so that $p$ and $q$ agree on all traces of
$D_0$.  Work in a common stabilized corner when needed.

Choose a countable local positive test family $s_i$, dense in linear
span in each of the basis ideals used in \Cref{seq:diagonal}.
Each test is cut-supported in both factors and has compact support
in $Y$. Choose an enlarged elementary test $\bar s_i$ and a fixed
$\kappa_i>0$ with $s_i\precsim(\bar s_i-\kappa_i)_+$, using the
positive cut in each factor. Hence
$p(s_i)\precsim(p(\bar s_i)-\kappa_i)_+$, and the same holds for $q$.
The model-independent upper sandwich, applied to this fixed cut of
$\bar s_i$, supplies a common compact-Cuntz-supported constant;
dominating it in the chosen containing open gives
$b_i\in Q(U_i)_+$ and a finite constant $L_i$ such that
\begin{equation}\label{un:controller-rank}
 d_\gamma(p(s_i))+d_\gamma(q(s_i))
 \leq L_i d_\gamma(b_i),\qquad\gamma\in T(D_0),
\end{equation}
for a finite $L_i$.  The constants and the comparisons are included
in $D_0$.  One can retain $b_i$ whenever later finite lists grow.

Fix a finite initial list and $\delta>0$.  For
$u\in\mathcal U(M_2(D_0)^{\sim})$, define the nonnegative continuous
functions
\[
 f_{u,i}(\gamma)
 =(\gamma\otimes\Tr_2)
 \bigl(\lvert u(p(s_i)\oplus0)u^*-(q(s_i)\oplus0)\rvert^2\bigr),
 \qquad \gamma\in K_{b_i}.
\]
Here $\Tr_2$ is unnormalized: the common zero padding leaves
\eqref{un:controller-rank} unchanged.  Continuity follows from
\Cref{un:trace-cap}.  We claim that zero lies
in the closed convex hull of the tuples $(f_{u,i})_i$ in
$\bigoplus_i C(K_{b_i},\R)$.  Otherwise a separating real functional
would be bounded below by a positive constant on all these tuples.
Represent its components by signed measures.  Because $f_{u,i}\geq0$,
replacing each measure by its positive part retains that lower bound.
By \Cref{un:barycenter}, the resulting functional on these functions
is a finite sum of positive multiples of evaluations at individual
weights in the corresponding caps.  \Cref{un:finite-traces} makes
all these evaluations simultaneously arbitrarily small, a
contradiction.

It follows that a finite convex combination satisfies
\begin{equation}\label{un:uniform-square}
 \sup_{\gamma\in K_{b_i}}
 \sum_{k=1}^r t_k f_{u_k,i}(\gamma)<\delta
 \qquad\text{for every tested }i.
\end{equation}
The coefficients may be chosen positive and rational: the finite
functions being averaged are bounded, so sufficiently small changes
of the coefficients preserve the strict inequality.  Write
$t_k=m_k/l$.  Before folding, work in $B_l=M_l(M_2(Q))$ with trace
$\tau\otimes\Tr_2\otimes\operatorname{tr}_l$: only this last,
additional averaging direction is normalized.  The two averaged maps
are diagonal copies of $\Ad(u_k)(p\oplus0)$ and of $q\oplus0$,
with $m_k$ occurrences of the $k$th block.  The original algebra $D_0$
has the fixed corner embedding
\[
 \iota_l:D_0\longrightarrow(B_l)_\infty,
 \qquad x\longmapsto(x\oplus0)\otimes1_l.
\]

We record explicitly the resulting coordinate estimate.  Lift the
finitely many unitaries to unitary multiplier sequences, and use
c.p.c. representatives of $p,q$.  Write $z_{i,n}$ for the difference
of their averaged values, after folding into $Q$.  For every fixed
$t>0$, we first prove the absolute estimate
\begin{equation}\label{un:cap-coordinate-rank}
 d_\tau((|z_{i,n}|-t)_+)
 \leq \frac{4\delta}{t^2}
 \qquad\bigl(d_\tau(b_i)\leq1\bigr)
\end{equation}
for every sufficiently large $n$, uniformly over this cap.  Suppose
that this absolute bound fails arbitrarily late, and choose
violating coordinates and weights $\tau_n$.  The trace identities
of \Cref{st:folding} pull these weights back, before folding, to
$\tau_n\otimes\Tr_2\otimes\operatorname{tr}_l$.  The subsequent multiplier
conjugations preserve these identities.  Positive cuts commute with
the coordinate folding maps, so regularization commutes with
them as well.  Let $\rho_l$ be the regularized trace before folding
and put $\gamma=\rho_l\circ\iota_l$.  It is a trace on the same
original $D_0$, and \Cref{seq:regularized} gives
\[
 d_\gamma(b_i)
 =\sup_{\eta>0}\lim_\omega d_{\tau_n}((b_i-\eta)_+)\leq1.
\]
Thus $\gamma\in K_{b_i}$.  Put
$e_{k,i}=u_k(p(s_i)\oplus0)u_k^*-(q(s_i)\oplus0)\in M_2(D_0)$.  If $z_i$ denotes the
unfolded diagonal difference, matrix-unit traciality and finite
additivity give
\begin{equation}\label{un:averaged-limit}
 \rho_l(z_i^*z_i)
 =\sum_k\frac{m_k}{l}(\gamma\otimes\Tr_2)(e_{k,i}^*e_{k,i})<\delta.
\end{equation}
Indeed, restriction to the first internal corner gives $\gamma$,
so the matrix extension on $M_2(D_0)$ is $\gamma\otimes\Tr_2$.
For every positive $a\in M_2(D_0)$,
$\rho_l(a\otimes e_{jj})=(\gamma\otimes\Tr_2)(a)/l$: the $l$
averaging corners are equivalent and sum to $a\otimes1_l$.  This uses only finite positive sums, so also holds
when a value is infinite.  It follows that
\eqref{un:averaged-limit} is valid without a bounded-weight
assumption.  On the other hand, functional calculus gives
\[
 \tau\bigl((|z_{i,n}|^2-t^2/2)_+\bigr)
 \geq \frac{t^2}{2}d_\tau((|z_{i,n}|-t)_+).
\]
Regularization and the violations of the absolute estimate therefore
give a square trace at least $2\delta$, contrary to
\eqref{un:averaged-limit}.  This proves
\eqref{un:cap-coordinate-rank}.  If $0<d_\tau(b_i)<\infty$, apply
that estimate to $\tau/d_\tau(b_i)$.  If $d_\tau(b_i)=0$, apply it
to every positive multiple of $\tau$, forcing the left side to be
zero.  For $d_\tau(b_i)=\infty$ the desired proportional bound is
automatic.  Thus, at the same sufficiently late coordinates,
\begin{equation}\label{un:coordinate-rank}
 d_\tau((|z_{i,n}|-t)_+)
 \leq \frac{4\delta}{t^2}d_\tau(b_i)
 \qquad\text{for all }\tau\in T(Q).
\end{equation}

At stage $j$ use the first $j$ local tests, the cut tolerances
$1,1/2,\ldots,1/j$, and choose
$\delta_j<2^{-j}/(4j^2)$.  Then advance far enough in that row for
\eqref{un:coordinate-rank}, the chosen algebraic tests, and their
norm bounds all to hold.  Pass through the rows slowly as in
\Cref{seq:diagonal}.  The averaging matrix size is fixed while
each row's uniform coordinate estimate is proved.  For a fixed local
test and any fixed cut
tolerance, this gives its $J(U_i)$ estimate against the fixed $b_i$.
Closure and the local density of the tests give
\eqref{un:local-difference} for all open ideals.

For completeness, these operations preserve the asserted relation
to the original maps.  At each coordinate, a diagonal sum of
conjugates of a map is conjugate, by the diagonal of the same
unitaries, to its tensor-first copy in the averaging matrix factor.
The coordinate folding can be chosen, using \Cref{st:folding}, to
carry that copy within the stage tolerance of the original on every
current input.  The first-corner clause of \Cref{st:folding} removes
the common internal zero padding at the same time.  Make that tolerance tend to zero during the same
diagonal choice.  Inner conjugacy preserves every coordinate weight,
and vanishing norm changes preserve regularized values by cuts, so
the models keep their data.  The argument applies to a designated
block of a direct sum, amplifying the spare by the same averaging
factor.  Its unused part can always be placed in a separate stable
corner before folding.
\end{proof}

\subsection{Ideal-related absorption and compact perturbations}

Tracial alignment has placed the difference of two models in each
local error ideal. To remove that difference in norm, we will dominate
its coefficients by a small stabilizer and apply stable uniqueness
on a hereditary support. We first state the absorption and compact
perturbation results used in that step.

We specify the version of the ideal-related machinery used below.
An $X$-structure on a coefficient algebra $H$ is an order-preserving
assignment $U\mapsto H(U)$ of closed ideals.  We do not assume that
this assignment preserves intersections.  Set
\[
 E_H=\ell^2(\N)\otimes H,\qquad
 \mathcal K(E_H)=H\otimes\K,\qquad
 \mathcal K(E_H)(U)=H(U)\otimes\K.
\]
A representation $\phi:S\to\mathcal L(E_H)$ is weakly equivariant
if
\begin{equation}\label{un:weak-equivariance}
 \langle\xi,\phi(a)\eta\rangle\in H(U)
 \quad(a\in S(U),\ \xi,\eta\in E_H).
\end{equation}
Equivalently, $\phi(S(U))E_H\subseteq E_HH(U)$.  Every adjointable
module operator preserves these submodules.  The diagonal coefficient
maps $a\mapsto\langle\xi,\phi(a)\xi\rangle$ are completely positive.
For a nuclear $S$ they and their induced quotient maps are nuclear,
since ideals and quotients of a nuclear $C^*$-algebra are nuclear
\cite[Corollary~4]{ChoiEffrosInjectivity}; other coefficients follow
by polarization.  Thus weak equivariance here
implies weak $X$-nuclearity, also for homotopies.

\Needspace{3\baselineskip}
\begin{theorem}[Ideal-related absorption]\label{un:absorption}
Let $S$ be separable, exact, and lower semicontinuous over $X$, and
let $H$ be sigma-unital.  Suppose
$\theta:S\to H\otimes\K$ is a nuclear $X$-full homomorphism, meaning
\[
 \Ideal_{H\otimes\K}(\theta(a))=H(U_a)\otimes\K
 \qquad(a\in S_+).
\]
Then its infinite repeat is weakly $X$-nuclear and absorbs every
weakly $X$-nuclear representation on the standard module, by unitary
equivalences with compact errors tending to zero on each input.
\end{theorem}

This is \cite[Theorem~11.14]{Gabe}, applied to $H\otimes\K$.
Neither separability of $H$ nor an absorption hypothesis on $H$ is
required.  The repeat is taken on the stabilized module.  In particular,
if a map $\theta_0:S\to H$ is first placed in the compact corner
$\theta_0\otimes e_{11}$, its stable repeat retains the other, zero
module slots.  After a permutation it has the form
$\operatorname{diag}(\theta_0,0,\theta_0,0,\ldots)$.
These zero slots remain available when absorbing a degenerate
representation.

For weakly equivariant representations $\phi,\psi:S\to
\mathcal L(E_H)$ with compact difference, write $[\phi,\psi]_X$ for
the class of the graded cycle
\begin{equation}\label{un:pair-cycle}
 \left(E_H\oplus E_H,
 \begin{pmatrix}\phi&0\\0&\psi\end{pmatrix},
 \begin{pmatrix}0&1\\1&0\end{pmatrix}\right).
\end{equation}
The two summands have opposite gradings.  The compact-difference
condition supplies the compact commutators in this cycle.  This
notation refers to the pair, without assigning separate classes to
the multiplier-valued maps.

\begin{proposition}[Multiplier stable uniqueness]\label{un:stable-uniqueness}
Suppose $S$ is separable, nuclear, and lower semicontinuous over $X$,
$H$ is sigma-unital, and $\theta:S\to H\otimes\K$ is an $X$-full
homomorphism.
Let $\phi,\psi:S\to\mathcal L(E_H)$ be weakly equivariant
representations with compact difference and $[\phi,\psi]_X=0$.
Writing $\theta^\infty$ for an infinite repeat, for every finite
$\mathcal F\subset S$ and $\eps>0$ there is a unitary
$v\in1+\mathcal K(E_H\oplus E_H)$ such that
\begin{equation}\label{un:proper-uniqueness}
 \|v(\phi\oplus\theta^\infty)(a)v^*
       -(\psi\oplus\theta^\infty)(a)\|<\eps,
 \qquad a\in\mathcal F.
\end{equation}
\end{proposition}

\begin{proof}
The multiplier-pair theorem \cite[Theorem~13.9]{Gabe} supplies a
weakly $X$-nuclear representation $\Theta$ and an identity-plus-compact
unitary $u$ comparing $\phi\oplus\Theta$ with $\psi\oplus\Theta$
to any prescribed finite-test tolerance $\delta$.  Here the pair
class may be read in the nuclear theory, since $S$ is nuclear.
Put $\sigma=\Theta\oplus\theta^\infty$.  By
\Cref{un:absorption}, there is a module unitary $V$ from the repeat
module to the module of $\sigma$ such that
\[
 \|V^*\sigma(a)V-\theta^\infty(a)\|<\eta,
 \qquad a\in\mathcal F.
\]
With $W_0=1\oplus V$, set
\[
 v=W_0^*(u\oplus1)W_0.
\]
It is a unitary, and $v-1$ is compact because conjugating a compact
module operator by a module unitary preserves compactness.  The
error in \eqref{un:proper-uniqueness} is at most $\delta+2\eta$.
Choose these parameters with $\delta+2\eta<\eps$.
This is the finite-test common-stabilizer transfer underlying
\cite[Lemmas~2.4 and~3.4]{DadarlatEilers}; no initial value of a
continuous unitary path is needed here.
\end{proof}

\begin{lemma}[The trivially labeled factor]\label{un:KK-nullity}
For the coarse $X$-structure on $S=C_0(Y)\otimes P$ and every
sigma-unital coefficient $X$-algebra $H$,
\[
 \KK(X;S,H)=0.
\]
In particular every weakly equivariant multiplier pair with compact
difference has zero class in \Cref{un:stable-uniqueness}.
\end{lemma}

\begin{proof}
Since $\W$ is simple and nuclear, every ideal of its spatial tensor
product with another $C^*$-algebra is a tensor-product ideal
\cite[Theorem~3.1]{GuptaJain}. Together with the transported labels of
\Cref{st:ideals}, this gives, after reordering the tensor factors, an
$X$-algebra $C=C_0(Y)\otimes A\otimes\K$ such that
\[
 S=\W\otimes C,\qquad S(U)=\W\otimes C(U).
\]
The identity of $\W$ has zero ordinary Kasparov class.  Exterior
tensoring its stabilized cycle relation and its homotopy with $C$
retains weak equivariance.  Explicitly, a typical coefficient of
the tensor representation, for $c\in C(U)$, is a finite sum
\[
 \sum_{i,j}\langle\xi_i,\pi(w)\xi_j\rangle
                  \otimes c_i^*c c_j\ \in\ \W\otimes C(U).
\]
The same formula applies throughout the homotopy.  It follows that
the identity cycle of $S$ is zero for these labels.  Taking its cycle
product with an $S$--$H$ cycle preserves this stabilized null relation.  The identity-product homotopy identifies the
same product with the original cycle, so that cycle is zero as well.

The cycle relation and mixed product statements used here are
\cite[Propositions~12.17 and~12.20 and Lemmas~12.18, 12.19,
and~12.23]{Gabe}.
They require the first algebra to be separable and the remaining
coefficient algebras to be sigma-unital.  Thus they apply to this
$H$ even when it is nonseparable.  Nuclearity of $S$ also makes the
cycles and homotopies weakly $X$-nuclear, as needed above.
\end{proof}

\begin{lemma}[Hereditary support and local multiplication]
\label{un:hereditary-support}
Let $I\subseteq C$ be a separable $C^*$-subalgebra with strictly
positive element $e$, and let $k\in C_+$.  Set
\[
 H=\Her_{M_2(C)}(e\oplus k).
\]
Then $H$ is sigma-unital, with no separability assertion required.
If $r\in C$ idealizes $I$, then $r\oplus0$ idealizes $H$.
If, in addition, $L$ is an ideal of $M_2(C)$ and
\[
 (rI)\oplus0\subseteq L,\qquad (Ir)\oplus0\subseteq L,
\]
then
\begin{equation}\label{un:support-bound}
 (r\oplus0)H+H(r\oplus0)\subseteq H\cap L.
\end{equation}
For $x\in H$, one has
$H\cap\Ideal_{M_2(C)}(x)=\Ideal_H(x)$.
\end{lemma}

\begin{proof}
The defining positive element is strictly positive in its hereditary
algebra, so $H$ is sigma-unital.  Since $r e^{1/2}$ and
$e^{1/2}r$ belong to $I$, multiplication of
$(e^{1/2}\oplus k^{1/2})M_2(C)(e^{1/2}\oplus k^{1/2})$
on either side by $r\oplus0$ remains in $H$.  These elements are
dense in $H$.  Under the additional hypothesis, the same products
also lie in $L$, proving \eqref{un:support-bound}.  The last identity
is the ideal correspondence for a hereditary subalgebra: an
approximate unit of $H$ compresses the ambient ideal generated by
$x$ to the ideal it generates in $H$.
\end{proof}

\begin{lemma}[Finite truncation of a proper conjugacy]
\label{un:finite-truncation}
Let $H$ be sigma-unital and let $p,q:S\to\mathcal M(H)$ be
representations with difference in $H$.  Let $\theta:S\to H$ be
a homomorphism, compact-valued on the module $H$, and form the diagonal
representations
\[
 A_0=p\oplus\theta\oplus\theta\oplus\cdots,
 \qquad B_0=q\oplus\theta\oplus\theta\oplus\cdots.
\]
Identical zero summands may be inserted in these diagonals.
If an identity-plus-compact unitary compares $A_0,B_0$ arbitrarily
well on finite sets, then, for every finite $\mathcal F$ and
$\eps>0$, finitely many copies suffice: for some $N,M\geq0$ there is a
unitary $v_N$ equal to the identity outside the initial summand,
the first $N$ repeat summands, and finitely many zero slots, with
\[
 \max_{a\in\mathcal F}
 \|v_N(p\oplus\theta^{\oplus N}\oplus0^{\oplus M})(a)v_N^*
       -(q\oplus\theta^{\oplus N}\oplus0^{\oplus M})(a)\|<\eps.
\]
All resulting errors are compact.  If the representations are
weakly equivariant, their compact errors on $S(U)$ belong to the
compacts over $H(U)$.
\end{lemma}

\begin{proof}
Choose $u=1+k$ with error less than $\eps/2$, and put
$R=\max(1,\max_{a\in\mathcal F}\|a\|)$.  If $E_N$ denotes the
specified finite-summand projection, choose
\[
 \|k-E_NkE_N\|<\eta,
 \qquad 0<\eta<\min\{1/4,\eps/(8R)\}.
\]
The operator $z_N=1+E_NkE_N$ is within $\eta$ of a unitary, so it
is invertible and the spectrum of $|z_N|$ is contained in
$[1-\eta,1+\eta]$.  Its polar unitary satisfies
\[
 v_N=z_N(z_N^*z_N)^{-1/2},\qquad
 \|v_N-z_N\|\leq\eta,
 \qquad\|v_N-u\|\leq2\eta.
\]
It is the identity off $E_N$ and differs from the identity by a
finite matrix over $H$.  Include successively the identical zero
slots as well, if present, so that $E_N\to1$ strongly.  The change
in the conjugacy error is at most $4R\eta$.  Since the two
representations are diagonal and agree on the omitted tail, the
remaining estimate is exactly the displayed finite-copy estimate.

For localization, a compact operator which is weakly equivariant
on an input in $S(U)$ has all matrix coefficients in $H(U)$,
hence belongs to $H(U)\otimes\K$.  All module unitaries preserve
the submodules $E_HH(U)$, so conjugation preserves this statement.
It applies to the compact differences here.
\end{proof}

\subsection{Removing the locally small difference}

\begin{proposition}[Uniqueness with a spare model]\label{un:spare}
Let $p,q:S\to D$ be equal-data models satisfying
$(p-q)(S(U))\subseteq J(U)$ for every open $U$.  Let $\chi$ be a
model for a strictly positive multiple of the same data.  Then
$p\oplus\chi$ and $q\oplus\chi$ are coordinate approximately
unitarily equivalent, with the matrix and trace conventions of
\Cref{st:folding}.
\end{proposition}

\begin{proof}
Put $C_1=C^*(p(S),q(S))$, and let $I$ be the ideal of $C_1$
generated by the differences $p(a)-q(a)$.  This is a separable
ideal.  For each open $U$,
\begin{equation}\label{un:difference-ideal}
 p(S(U))I+q(S(U))I+Ip(S(U))+Iq(S(U))\subseteq J(U).
\end{equation}
Indeed, modulo $J(U)$, the two maps agree on $S(U)$.  Their common
image is an ideal in the algebra generated by both images, because
multiplication from either side by $p(a)$ or $q(a)$ corresponds to
multiplication inside the source ideal.  It annihilates each
difference: for $s\in S(U)$,
\[
 p(s)(p(a)-q(a))=p(sa)-q(sa)=0\pmod{J(U)},
\]
and the other side is identical.  It therefore annihilates the
ideal generated by the differences, proving
\eqref{un:difference-ideal}.

Use the existing-model version of \Cref{sm:stabilizers}, with
the spare $\chi$, to obtain a locally small homomorphism $\theta$.
Impose the coefficient dominations
\begin{equation}\label{un:theta-domination}
 p(S(U_x))I+q(S(U_x))I+Ip(S(U_x))+Iq(S(U_x))
 \subseteq\Ideal(\theta(x)),\qquad x\in S_+.
\end{equation}
Apply \Cref{sm:coefficient-domination} to the four bounded
bilinear maps $(a,b)\mapsto p(a)b$, $q(a)b$, $bp(a)$, and $bq(a)$,
with $b\in I$.  Equation~\eqref{un:difference-ideal} gives their
local hypotheses, and that corollary supplies the displayed inclusion
for every positive $x$.  The first version also ensures
\begin{equation}\label{un:theta-range-fullness}
 \theta(S(U_x))\subseteq\Ideal(\theta(x)).
\end{equation}
Moreover $\theta$ is represented by the same coordinate maps as
$\chi$, tensored with rational UHF projections of trace tending
to zero.

Choose strictly positive elements $e\in I$ and $k\in\theta(S)$,
and form
\[
 H=\Her_{M_2(D)}(e\oplus k).
\]
The maps $p,q$, placed in the first corner and zero in the second,
idealize $H$, by \Cref{un:hereditary-support}.  View them as
representations on the module $H$, and then stabilize that module.
Equip $H$ with the order-preserving ideals
\begin{equation}\label{un:H-action}
 H(U)=\Ideal_H(0\oplus\theta(S(U))).
\end{equation}
No intersection property of this assignment is needed.  From
\eqref{un:theta-domination}, \Cref{un:hereditary-support}, and its
ideal-intersection identity, values of the first-corner maps on
$S(U_x)$ multiply $H$ into
$\Ideal_H(0\oplus\theta(x))$.  Thus $p,q$ are weakly equivariant
for \eqref{un:H-action}.  By \eqref{un:theta-range-fullness},
$0\oplus\theta(x)$ generates $H(U_x)$, so this compact-valued map
is $X$-full.  The algebra $H$ is sigma-unital; it is not required
to be separable.

The multiplier pair agrees modulo $H$, since its differences
belong to $I\oplus0\subseteq H$.  Pad the module $H$ by zero
copies to the standard module; the resulting pair class
\eqref{un:pair-cycle} is zero by \Cref{un:KK-nullity}.
The coarse source action is lower semicontinuous, since the
projection $Y\times X\to X$ is open.  For
\Cref{un:stable-uniqueness}, place $0\oplus\theta$ in the first
compact corner of the standard module.  Its infinite repeat
is, after reindexing, the diagonal of copies of $\theta$ on $H$
with identical zero slots.  Thus \Cref{un:finite-truncation}
applies on the original $H$ summands, yielding finite matrices
over $H$, rather than retaining an additional untruncated compact
direction inside each slot.  For any finite $\mathcal F$ and
$\eps>0$ this gives a finite $N$ and an identity-plus-finite-matrix
unitary comparing
\begin{equation}\label{un:finite-spare-comparison}
 p\oplus\theta^{\oplus N}
 \quad\text{and}\quad
 q\oplus\theta^{\oplus N}
\end{equation}
within $\eps$, with zero slots permitted.

These estimates hold for the actual ambient elements, as follows.
If $v=1+c$ is the obtained finite unitary and $A,B$ denote the
two finite diagonal maps, then
\begin{equation}\label{un:ambient-error}
 vA(a)v^*-B(a)
 =(A(a)-B(a))+cA(a)+A(a)c^*+cA(a)c^*.
\end{equation}
Every term belongs to a finite matrix over $H$, because the maps
idealize $H$ and their difference belongs to it.  The multiplier
norm of this element is its intrinsic $H$-norm and hence its
ambient norm in the matrix algebra over $D$.  This argument does
not assume that restriction of the entire ambient idealizer to
$\mathcal M(H)$ is faithful.  Only its isometric restriction to
$H$ is used.

There is also no lifting obstruction for the finite unitary.
Write it as $1+c$ in the unitization of a finite matrix algebra
over $D$, and choose bounded representatives $c_n$ over $Q$.
The two unitary defects of $1+c_n$ tend to zero.  For all
sufficiently large $n$ its polar part
\[
 v_n=(1+c_n)\bigl((1+c_n)^*(1+c_n)\bigr)^{-1/2}
\]
is a unitary in the corresponding matrix multiplier algebra and
differs from $1+c_n$ by a sequence tending to zero.  Thus a strict
quotient estimate, first obtained with a smaller tolerance, holds
at every sufficiently late coordinate.

Finally fix this $N$ before advancing the coordinates.  Since the
UHF trace of the projection defining $\theta_n$ tends to zero,
$N$ disjoint equivalent copies fit inside the positive spare
fraction eventually.  Using exactly the same $\chi_n$ in those
copies identifies them with \eqref{un:finite-spare-comparison}.
Move the unused spare to a separate stable corner and leave it
unchanged.  Consequently the finite conjugacy need not preserve
an internal complement of its original matrix slots.  Apply the
coordinate folding and the multiplier permutations from
\Cref{st:folding}, on all current inputs and all these blocks.
This proves eventual finite-test conjugacy of $p\oplus\chi$ and
$q\oplus\chi$.  On increasing finite sets with tolerances tending
to zero, choose each coordinate beyond the corresponding
threshold.  The resulting single sequence of multiplier unitaries
proves the proposition.
\end{proof}

\begin{proof}[Proof of \Cref{un:uniqueness}]
Split the tensor-first copy of each map along two equivalent UHF
projections of trace $1/2$.  Denote its two halves by $p_0,p_1$
and $q_0,q_1$.  Both halves have the original model data multiplied
by $1/2$, and their sums are coordinate approximately unitarily
equivalent to the original maps by \Cref{st:folding}.

First compare $p_0\oplus p_1$ with $q_0\oplus p_1$, retaining
$p_1$ as a common spare.  Apply \Cref{un:alignment} in the changing
block to put its difference in every appropriate $J(U)$, carrying
the spare through the same diagonal amplifications.  Then
\Cref{un:spare} gives the required conjugacy.  The coordinate
conjugacies relating the aligned blocks to the originals may be
undone, with their vanishing errors, so the conclusion applies
to the displayed original half-sums.  Next compare
$q_0\oplus p_1$ with $q_0\oplus q_1$, using $q_0$ as the common
spare and interchanging the two matrix positions.  The same
argument applies.  Composing the two coordinate conjugacies and
the folding adjustments proves the assertion on every finite
test set and tolerance.  A diagonal choice on a dense increasing
family gives the stated sequence $u_n$.

Finally, any two c.p.c. lifts of a fixed model differ in norm by a
sequence tending to zero on each input.  The conclusion therefore
holds for the lifts specified in the statement.
\end{proof}
